\par\Needspace{6\baselineskip}\medskip
\hypertarget{native-151-article109}{}\pdfbookmark[2]{109.}{article-109}
\begin{center}109.\end{center}
Verwenden wir für die Hyperbel die Zeichen $\alpha$, $C$ und $c$ in derselben Bedeutung wie in Artikel 99, so ergibt sich aus den Gleichungen VIII und IX in Artikel 21
\[
\begin{aligned}
r'\cos v'-r\cos v&=-\tfrac12\left(c-\frac1c\right)\left(C-\frac1C\right)\alpha\\[.7em]
r'\sin v'-r\sin v&=\tfrac12\left(c-\frac1c\right)\left(C+\frac1C\right)\alpha\BareRoot{(ee-1)},
\end{aligned}
\]
\noindent und somit
\[
\rho=\tfrac12\alpha\left(c-\frac1c\right)\BareRoot{\left(ee\left(C+\frac1C\right)^2-4\right)}.
\]
\noindent\SourceLeaf{152}Sei $\gamma$ die durch die Gleichung $\gamma+\dfrac1\gamma=e\left(C+\dfrac1C\right)$ bestimmte Größe. Da dieser Gleichung offenbar \emph{zwei} zueinander reziproke Werte genügen, wählen wir denjenigen, der größer als 1 ist. Damit ergibt sich
\[
\rho=\tfrac12\alpha\left(c-\frac1c\right)\left(\gamma-\frac1\gamma\right)\mathpunct{\cdot}
\]
\noindent Ferner gilt
\[
\begin{aligned}
r+r'&=\tfrac12\alpha\left(e\left(c+\frac1c\right)\left(C+\frac1C\right)-4\right)\\[.7em]
&=\tfrac12\alpha\left(\left(c+\frac1c\right)\left(\gamma+\frac1\gamma\right)-4\right),
\end{aligned}
\]
\noindent und somit
\[
\begin{aligned}
r+r'+\rho&=\alpha\left(\BareRoot{c\gamma}-\BareRoot{\frac1{c\gamma}}\right)^2\\[.7em]
r+r'-\rho&=\alpha\left(\BareRoot{\frac\gamma c}-\BareRoot{\frac c\gamma}\right)^2\mathpunct{\cdot}
\end{aligned}
\]
\noindent Setzen wir also
\[
\BareRoot{\frac{r+r'+\rho}{4\alpha}}=m,\qquad\BareRoot{\frac{r+r'-\rho}{4\alpha}}=n,
\]
\noindent so gilt notwendig $\BareRoot{c\gamma}-\BareRoot{\dfrac1{c\gamma}}=2m$. Um zu entscheiden, ob $\BareRoot{\dfrac\gamma c}-\BareRoot{\dfrac c\gamma}$ den Wert $=+2n$ oder $=-2n$ annimmt, müssen wir feststellen, ob $\gamma$ größer oder kleiner als $c$ ist. Aus Gleichung 8 in Artikel 99 folgt leicht, dass der erste Fall immer dann eintritt, wenn $2f$ unter $180^0$ liegt, und der zweite immer dann, wenn $2f$ über $180^0$ liegt. Schließlich ergibt sich aus demselben Artikel
\[
\begin{aligned}
\frac{kt}{\alpha^{\frac32}}&=\tfrac12\left(\gamma+\frac1\gamma\right)\left(c-\frac1c\right)-2\log c\\[.8em]
&=\tfrac12\left(c\gamma-\frac1{c\gamma}\right)-\tfrac12\left(\frac\gamma c-\frac c\gamma\right)-\log c\gamma+\log\frac\gamma c\\[.8em]
&=2m\BareRoot{(1+mm)}\mp2n\BareRoot{(1+nn)}\\[.5em]
&\qquad-2\log\bigl(\BareRoot{(1+mm)}+m\bigr)\pm2\log\bigl(\BareRoot{(1+nn)}+n\bigr)
\end{aligned}
\]
\noindent wobei die unteren Vorzeichen stets für den Fall $2f>180^0$ gelten. Nun lässt sich
\[
\log\bigl(\BareRoot{(1+mm)}+m\bigr)
\]
\noindent leicht in die folgende Reihe entwickeln
\[
m-\tfrac13\cdot\tfrac12m^3+\tfrac15\cdot\frac{1\mathpunct{.}3}{2\mathpunct{.}4}m^5-\tfrac17\cdot\frac{1\mathpunct{.}3\mathpunct{.}5}{2\mathpunct{.}4\mathpunct{.}6}m^7+\text{usw.}
\]
\noindent Das folgt unmittelbar aus $\mathrm d\log\bigl(\BareRoot{(1+mm)}+m\bigr)=\dfrac{\mathrm dm}{\BareRoot{(1+mm)}}\mathpunct{\cdot}$ Damit erhalten wir
\[
\begin{aligned}
&2m\BareRoot{(1+mm)}-2\log\bigl(\BareRoot{(1+mm)}+m\bigr)\\[.6em]
&\qquad=4\left(\tfrac13m^3-\tfrac15\cdot\tfrac12m^5+\tfrac17\cdot\frac{1\mathpunct{.}3}{2\mathpunct{.}4}m^7-\text{usw.}\right)
\end{aligned}
\]
\noindent\SourceLeaf{153}und ebenso eine ganz entsprechende Formel, wenn man $m$ mit $n$ vertauscht. Setzen wir schließlich
\[
\begin{aligned}
T={}&\tfrac16(r+r'-\rho)^{\frac32}-\tfrac1{80}\cdot\frac1\alpha(r+r'-\rho)^{\frac52}\\[.6em]
&+\tfrac3{1792}\cdot\frac1{\alpha\alpha}(r+r'-\rho)^{\frac72}-\tfrac5{18432}\cdot\frac1{\alpha^3}(r+r'-\rho)^{\frac92}+\text{usw.}
\end{aligned}
\]
\[
\begin{aligned}
U={}&\tfrac16(r+r'+\rho)^{\frac32}-\tfrac1{80}\cdot\frac1\alpha(r+r'+\rho)^{\frac52}\\[.6em]
&+\tfrac3{1792}\cdot\frac1{\alpha\alpha}(r+r'+\rho)^{\frac72}-\tfrac5{18432}\cdot\frac1{\alpha^3}(r+r'+\rho)^{\frac92}+\text{usw.,}
\end{aligned}
\]
\noindent so folgt $kt=U\mp T$. Diese Ausdrücke stimmen vollständig mit den in Artikel 107 angegebenen überein, wenn man dort $a$ durch $-\alpha$ ersetzt.
\par
Sowohl für die Ellipse als auch für die Hyperbel sind diese Reihen in der Praxis besonders bequem, wenn $a$ oder $\alpha$ einen sehr großen Wert annimmt, das heißt, wenn der Kegelschnitt einer Parabel sehr nahekommt. In einem solchen Fall ließen sie sich auch zur Lösung der oben behandelten Aufgabe (Artikel 85--105) verwenden. Da sie aber nach unserem Urteil selbst dann nicht die Kürze der oben angegebenen Lösung bieten, gehen wir auf dieses Verfahren nicht näher ein.
\par\bigskip\begin{center}\rule{.23\linewidth}{.35pt}\end{center}
